\newpage
\section{Gateway module software}

The gateway module is the most important part of this project as it interfaces with
the car and is resposible for the entire system's power management. The module receives
and parses VAN frames, organizes the data and then forwards it to the infotainment computer.
It also does the reverse. It can receive data from the infotainment computer and transmit the
appropriate VAN frame on the bus. The power management is also an important part of the software,
as excessive idle power draw will drain the car's battery fairly quickly.

\subsection{Framework and approach}

For the software, we are using the ESP-IDF framework from Espressif. The framework is based on
the FreeRTOS Operating System. Due to the nature of the OS and the software environment, the module
software was written with an event-based approach.

\bigskip
\noindent
We use the following events as a base for the main loops/tasks:
\begin{itemize}
  \item Receiving VAN frames
  \item Receiving UART data from the computer
  \item Car state change
  \item Timer events
  \item Main task
\end{itemize}

\noindent
Since the ESP32 has two cores, we delegated one core to exclusively handle receiving and parsing VAN frames,
since they are very frequent and we don't want to miss any.

\subsection{Receiving VAN frames}

To receive VAN frames we have to manually read and time the RX pin states. Luckily, the ESP32
has a hardware peripheral that does just that: the RMT peripheral. It can monitor GPIO pin states
and also measure how long it stayed in that state. Once we set the thresholds for the minimum and maximum
signal lengths we can have the peripheral trigger an interrupt every time it sees an \textbf{EOF} sequence.
\begin{listing}[H]
  \begin{minted}{C}
  rmt_rx_channel_config_t rx_chan_config = {0};
  
  ...
  
  timeslot_interval_ns = 8 * 1000; // At 125 kTS/s
  
  // EOF is 10 TS long
  rx_recv_config.signal_range_max_ns = 10 * timeslot_interval_ns;
  rx_recv_config.signal_range_min_ns = timeslot_interval_ns / 2;
  
  ...
  \end{minted}
  \caption{RMT timing initialization}
\end{listing}

\noindent
Once the \textbf{EOF} sequence is found, the RMT peripheral driver triggers an interrupt. Inside that interrupt,
we put a new event in the VAN Receiver task queue and restart the peripheral. The VAN Receiver task takes the
received RMT data from its queue and parses it. The data is an array of structs whose defition can be found in
Listing \ref{lst:rmt_symbol_word_t}. We reconstruct the VAN frame bit by bit from that array, taking into account
the Manchester bits. After checking the checksum, we separate the identifier and data, put it in a new buffer and
post it to the queue of the Main task for further processing. The receiving logic was inspired by
Peter Pinter's arduino library \cite{esp32_rmt_van} and adapted to suit our firmware architecture.

\begin{listing}[H]
  \begin{minted}{C}
/**
 * @brief The layout of RMT symbol stored in memory, 
 *        which is decided by the hardware design
 */
typedef union {
    struct {
        uint16_t duration0 : 15; /*!< Duration of level0 */
        uint16_t level0 : 1;     /*!< Level of the first part */
        uint16_t duration1 : 15; /*!< Duration of level1 */
        uint16_t level1 : 1;     /*!< Level of the second part */
    };
    uint32_t val; /*!< Equivalent unsigned value for the RMT symbol */
} rmt_symbol_word_t;
  \end{minted}
  \caption{Definiton of the RMT data symbol in ESP-IDF}
  \label{lst:rmt_symbol_word_t}
\end{listing}

\subsection{Parsing VAN frames}

Once the VAN receiver task posts a new event, the Main task then calls the packet parser which checks
the identifier and, if it's recognized, parses the data into global buffers and posts appropriate events
to the Main queue depending on the packet being parsed. See listing \ref{lst:example-van-frame} for an
example VAN frame struct. The vast majority of known decoded VAN frames can be found on Peter Pinter's
website \cite{vanbus_packets}.

\begin{listing}[H]
  \begin{minted}{C}
/**
 * @brief VAN rpm and speed data struct. Iden 0x824
 *
 */
struct psa_van_rpm_speed_struct
{
    uint16_t rpm;           // RPM in big endian. Divide by 8 to get actual rpm;
    uint16_t speed;         // Speed in big endian. Divide by 100 to get actual speed;
    uint16_t distance;      // Distance covered. Increments with each passed decimeter
    uint8_t packet_changed; // Increments each time information significantly changes
} __attribute__((packed));
  \end{minted}
  \caption{Example VAN frame struct.}
  \label{lst:example-van-frame}
\end{listing}

\subsection{Sending VAN frames}

We use the TSS463C to send VAN frames to the car. The TSS is effectively additional 128 bytes of RAM
that we have to manage, since we have to manage all the memory used for sending and receiving data.
For now, the TSS is only used to send CD Changer packets so we can emulate it, and to send trip reset
requests. The CD Changer messages are immediate reply messages (see \ref{subsec:message-types}), and the
trip reset is a normal frame with an ACK request. The CD Changer packets are on a timer and sent every
second, otherwise the onboard radio will think the CD Changer is absent and disable it.

\noindent
The TSS communicates over SPI and is a bit cumbersome to use, so a custom driver was written to make it easier to use. Additionally,
there is a seperate task that monitors the TSS message states and queues them to be sent.

\subsection{UART communication}

The infotainment computer communicates with the ESP32 over UART. The reason we use UART is its simplicity of use.
The ESP mostly sends data to the computer, but can also receive requests from the computer.

\noindent
The uart code lives in its own task that receives events from other tasks and can also send events to the main task.
When the Main loop sees a change to the global data buffers (a VAN frame arrived, battery voltage changed, etc\dots)
it posts an appropriate event to the uart task with the data to send to the computer. The protocol and data structs
are custom and inspired by the MSP protocol used in Flight Controller firmware on FPV drones.

\subsection{Event handling and tasks}

The firmware uses several tasks that run as their own loops and wait for events from multiple queues. There are a total
of 4 global queues in the firmware, one for each major component:
\begin{itemize}
  \item Main Queue - Read by the Main task. All other tasks post their events to the Main task queue.
        The Main task then determines what to do and where to forward that event next.
  \item TSS Queue - Read by the TSS task. Used to queue packets for transmission by the TSS.
  \item UART Queue - Read by the UART task. Used to send packets to the Infotainment computer.
  \item VAN Queue (Unused) - Read by the VAN RMT task. Left for future use.
\end{itemize}

\noindent
Additionally, there is an internal queue in the VAN RMT task. This queue is used exclusively between the RMT
peripheral's interrupt and the VAN RMT task.

\bigskip
\noindent
The events are defined as a struct (Listing \ref{lst:mive_global_event}) with the event identifier
(Listing \ref{lst:enum-mive_global_event_t}) and a pointer to some data. The data depends on the
event being received.

\subsection{Handling power states}

It is important we manage the system's power consumption. To make that easier, we have defined
several `car states'. See table \ref{tbl:car-states} for a list of those states, and the systems
behaviour. One important component mentioned in that table is ESP32's deep sleep mode. When placed
in deep sleep, the ESP is practically turned off. The only thing running on the chip is a special
section of RAM, the RTC peripheral and, if needed, the ULP coprocessor. For our usecase, we use the
same ADC pin we use for battery voltage measurement to wake up the ESP. Due to the voltage ranges
of the pin, we can't use simple GPIO wake up, but must use the ULP coprocessor to measure voltage
and wake the ESP up when it increases past a certain threshold.

\subsection{The ULP coprocessor}

The ESP32 has a very low power coprocessor onboard. It has a very simple architecture and uses very little
power. It is ideal to use for ADC measurement and waking up the ESP. It cannot be programmed directly in C,
but has two other ways of writing code. The first one is writing the assembly code in a seperate file and
building it into the firmware binary. The other one is using predefined C macros and writing the program
in an array inside the C program and loading it that way. We opted to use the second option as it's
way simpler to use. See Listing \ref{lst:ulp-deep-sleep} for the ULP code used in deep sleep.

\begin{table}[H]
  \begin{tabularx}{\textwidth} { 
  | >{\raggedright\arraybackslash}X 
  | >{\raggedright\arraybackslash}X 
  | >{\raggedright\arraybackslash}X | }
    \hline
    \textbf{Description} & \textbf{Computer state} & \textbf{ESP32 state} \\
    \hline
    \hline
    Car in sleep mode (switched 12V rail is off) & Turned Off (no power) & Deep sleep \\
    \hline
    Car off and locked, not in sleep mode & Turned Off & Running \\
    \hline
    Car off, not locked, not in sleep mode & If coming from states above, turned off, otherwise display off & Running \\
    \hline
    Car off, radio turned on & Turned on, display on & Running \\
    \hline
    Car off, key in ACC & Turned on, display on & Running \\
    \hline
    Car off, key in IGN & Turned on, display on & Running \\
    \hline
    Car cranking & Turned on, display off & Running \\
    \hline
    Engine running & Turned on, display on & Running \\
    \hline
  \end{tabularx}
  \caption{Car state table}
  \label{tbl:car-states}
\end{table}

\begin{listing}[H]
  \begin{minted}{C}
struct mive_global_event
{
    enum mive_global_event_t event;
    void* ev_data;
};
  \end{minted}
  \caption{Event struct definition}
  \label{lst:mive_global_event}
\end{listing}

\begin{listing}[H]
  \begin{minted}{C}
enum mive_global_event_t
{
    MIVE_EVENT_NONE = 0,
    MIVE_EVENT_GLOBAL_SLEEP,
    MIVE_EVENT_RMT_NEW_VAN_FRAME,
    MIVE_EVENT_TSS_RESET,
    MIVE_EVENT_TSS_ACTIVATE,
    MIVE_EVENT_TSS_IDLE,
    MIVE_EVENT_TSS_SLEEP,
    MIVE_EVENT_TSS_WRITE_FRAME,
    MIVE_EVENT_TSS_MONITOR_CHANNEL,

    MIVE_EVENT_VAN_UPDATE_AUDIO_MENU,
    MIVE_EVENT_VAN_NEXT_AUDIO_MENU_ITEM,
    MIVE_EVENT_VAN_CLOSE_AUDIO_MENU,

    MIVE_EVENT_UART_NEW_DATA, // New VAN Data to send
    MIVE_EVENT_UART_RECEIVE, // Received data from UART
    MIVE_EVENT_UART_SEND, // Send data to UART

    MIVE_EVENT_TIMER_UART_SEND,
    MIVE_EVENT_TIMER_1MS,
    MIVE_EVENT_TIMER_1S,

    MIVE_EVENT_ADC,
    MIVE_EVENT_STATE_CHANGE,
    MIVE_EVENT_POWER,
};
  \end{minted}
  \caption{List of some events present in the firmware}
  \label{lst:enum-mive_global_event_t}
\end{listing}

\begin{listing}[H]
  \begin{minted}{C}
const ulp_insn_t program[] = {
    I_MOVI(R0, 0),                // Set reg. R0 to initial 0
    I_MOVI(R2, 0),                // Set reg. R2 to initial 0
    M_LABEL(1),                   // LABEL 1 - Start of measurement
    I_ADDI(R0, R0, 1),            // Increment cycle counter (reg. R0)
    I_ADC(R1, PSA_ADC_UNIT, PSA_ADC_CHANNEL),            // Read ADC value to R1
    I_ADDR(R2, R2, R1),           // Add ADC value from R1 to R2
    M_BL(1, 4),                   // If cycle counter is less than 4, go to LABEL 1
    I_RSHI(R0, R2, 2),            // Divide accumulated ADC value in R2 by 4 and save it to R0
    M_BGE(2, high_adc_treshold),  // If average ADC value from reg. 
                                  // R0 is higher or equal than high_adc_treshold, go to LABEL 2
    I_HALT(),                     // Halt the coprocessor, it will be woken up by the timer again
    M_LABEL(2),                   // LABEL 2
    I_WAKE(),                     // Wake up ESP32
    I_END(),                      // Stop ULP program timer
    I_HALT(),                     // Halt the coprocessor
};

size_t size = sizeof(program)/sizeof(ulp_insn_t);
ESP_ERROR_CHECK(ulp_process_macros_and_load(0, program, &size));
ulp_set_wakeup_period(0, 20000); // ULP wakeup timer
err = ulp_run(0);
  \end{minted}
  \caption{ULP code that runs in deep sleep.}
  \label{lst:ulp-deep-sleep}
\end{listing}

\begin{listing}[H]
  \begin{minted}{C}
int ret = xQueueReceive(main_queue, &event, pdMS_TO_TICKS(1000));

if(ret == pdPASS)
{
    switch (event.event)
    {
    case MIVE_EVENT_RMT_NEW_VAN_FRAME:
        van_packet = (mive_van_packet_t*)event.ev_data;

        ret = psa_parse_van_packet(
            van_packet->iden,
            van_packet->packet_size,
            van_packet->packet,
            &global_libpsa_buffers);
        
        ...
        break;
    }
}
  \end{minted}
  \caption{Example event handling}
\end{listing}